\subsection{Post hoc capacity check: a small mapping can be memorized}

The low primary scores motivated a separate capacity diagnostic after the main
matrix had completed.
A newly initialized main-architecture bridge was fitted to 16 generated questions
with data seed 20261110, batch size 8 and 100 epochs.
For this diagnostic only, validation was an exact copy of those training items,
and checkpoint selection minimized resubstitution loss.
The dataset metadata explicitly marks the intended overlap.
This condition measures memorization capacity, not held-out performance, and
changes both dataset size and the number of exposures per question.

The fitted bridge answers all 16 training questions correctly.
When their source prefixes are shuffled, only 3/16 answers match the original
target references, but all 16 outputs equal the answer associated with the
\emph{supplied source} question.
The three remaining target matches consist of one fixed point in the permutation
and two coincident reference answers.
Thus the interface can memorize a small source-dependent mapping; its poor main
result is not evidence that the computational path is incapable of transmitting
any learned correspondence.
Together with the main fit's 8.7\% training accuracy, the diagnostic points to
difficulty fitting the larger problem under the selected representation and
optimization budget, rather than a failure confined to unseen questions.
It does not identify which architectural or training change would resolve that
difficulty.

The diagnostic has its own configuration, data, histories, source-shuffle mapping
and raw predictions under \path{overfit_diagnostic}; its resubstitution scores are
excluded from Table~\ref{tab:accuracy} and the primary accuracy figure.
